\bookchapter{الفصل 16: نظام الحسابات والخدمات السحابية}

\chapterintro{لي صديق، سُرق هاتفه العام الماضي. أول ما فكّر فيه لم يكن “كم ثمن الهاتف”، بل “هل يمكنني استرجاع صوري؟”. للأسف، لم يكن قد فعّل النسخ الاحتياطي السحابي. الهاتف شيء خارجي، لكن البيانات داخل حسابك، إذا ضاعت حقاً فلن تعود أبداً. في هذا الفصل سنتحدث عن كيفية تأمين حياتك الرقمية بشكل حقيقي.}

\booksection{أولاً: حساب هاتفك، ماذا يدير فعلاً؟}

كثيرون يظنون أن الحساب فقط لتنزيل التطبيقات. في الواقع، ما يربطه حساب الهاتف أكثر بكثير مما تتخيل:

\begin{itemize}
  \item \textbf{جهات الاتصال}: أرقام هواتف جمعتها لسنوات أو حتى لعشرات السنين.
  \item \textbf{الصور والفيديوهات}: من ولادة الطفل إلى حفل التخرج، من الموعد الأول إلى رحلة الزفاف.
  \item \textbf{الملاحظات}: أفكارك ومذكراتك وكلمات مرورك المسجلة على عجل (نعم، كثيرون يخزنون كلمات المرور في الملاحظات).
  \item \textbf{سجل مشتريات التطبيقات}: التطبيقات المدفوعة والمشتريات داخل التطبيقات التي أنفقت عليها مئات اليوانات.
  \item \textbf{العثور على الهاتف}: بعد فقدان الهاتف، هذه الأداة الوحيدة التي تساعدك على تحديد الموقع وقفل الجهاز ومسح البيانات عن بُعد.
\end{itemize}

إذا كنت تستخدم هاتف Apple، فكل ما سبق مرتبط بـ Apple ID. وإذا كنت تستخدم أندرويد، فهو مرتبط بحساب Huawei/Xiaomi/OPPO/vivo، وحساب Google الخاص بك. هذه الحسابات هي “سجل العائلة” لحياتك الرقمية، وإذا فقدت السجل، فإن الكثير من الأشياء تصبح غير قابلة للإثبات.

\booksection{ثانياً: ثلاثة أمور يجب فعلها، الآن فوراً}

\begin{enumerate}
  \item \textbf{تفعيل المصادقة الثنائية}: سواء كنت تستخدم Apple أو أندرويد، أول شيء هو الذهاب إلى إعدادات الحساب وتفعيل "المصادقة الثنائية" أو "التحقق بخطوتين". هذا الأمر مثل إضافة قفل إضافي على بابك — كلمة المرور وحدها لا تكفي، بل يجب أن يتلقى هاتفك رمز تحقق. حتى لو خمّن أحدهم كلمة مرورك عبر الإنترنت، بدون هاتفك لن يتمكن من الدخول.
  \item \textbf{ربط طريقة استرجاع موثوقة}: المصادقة الثنائية آمنة جداً، لكن لها مخاطرة — ماذا لو فقدت هاتفك نفسه؟ لذلك سيطلب منك النظام إعداد "مفتاح الاسترجاع" أو "بريد/هاتف احتياطي". هذا هو طريقك الوحيد لاستعادة الحساب في الطوارئ، يُفضّل كتابته على ورقة، ولا تخزنه في الهاتف.
  \item \textbf{التأكد من تفعيل "العثور على هاتفي"}: هذه الميزة تساعدك بعد فقدان الهاتف على تحديد الموقع، والقفل، وتشغيل الصوت، بل ومسح جميع البيانات عن بُعد. إنها خط الدفاع الأخير. كثيرون يتخطون هذه الخطوة عند شراء الهاتف، وعندما يحتاجونها لاحقاً يكتشفون أنها لم تُفعّل، فيكون الوقت قد فات.
\end{enumerate}

\booksection{ثالثاً: iCloud: “الهواء” لمستخدمي Apple}

من يستخدم هاتف Apple، سيواجه عاجلاً أم آجلاً سؤالاً وجودياً: ماذا أفعل عندما تمتلئ مساحة iCloud؟

تمنحك Apple مجاناً 5 غيغابايت، ولنكون صادقين، هذه الـ 5 غيغابايت لا تكفي لنسخة احتياطية واحدة. لذلك سيحصل معظم الناس في النهاية على تلك النافذة المزعجة: “مساحة تخزين iCloud ممتلئة”. لديك ثلاثة خيارات:

\begin{itemize}
  \item \textbf{التنظيف}: اذهب إلى الإعدادات اسمك iCloud إدارة التخزين، لترى من يحتل مساحتك. عموماً، الصور والنسخ الاحتياطية هي الأكبر. يمكنك البدء بحذف النسخ الاحتياطية القديمة غير الضرورية.
  \item \textbf{الترقية}: 50 غيغابايت بـ 6 يوانات شهرياً، 200 غيغابايت بـ 21 يواناً شهرياً. أرى شخصياً أن 50 غيغابايت تكفي لمعظم الناس، تُستخدم أساساً للنسخ الاحتياطي وتخزين صور مصغرة.
  \item \textbf{حل بديل}: إذا لم ترغب في الدفع، يمكنك تصدير الصور دورياً إلى الحاسوب، ثم إيقاف مزامنة صور iCloud. لكنك ستفقد بذلك راحة النسخ الاحتياطي التلقائي للصور.
\end{itemize}

تذكير مهم: مكتبة صور iCloud ليست “قرصاً سحابياً” — بل هي “مزامنة”. إذا حذفت صورة من هاتفك، تُحذف من iCloud أيضاً. إذا كنت تريد “النسخ الاحتياطي إلى السحابة ثم إفراغ الهاتف”، فصور iCloud ليست ما تبحث عنه، بل تحتاج إلى iCloud Backup، أو قرص سحابي من طرف ثالث.

\begin{tcbinfo}{iCloud+: ثلاثية الخصوصية للمستخدمين المدفوعين فقط}
بمجرد الترقية إلى أي فئة مدفوعة من iCloud (50GB/200GB/2TB وغيرها)، تصبح تلقائياً\textbf{iCloud+}مستخدم، وتفتح ثلاث وظائف حماية الخصوصية:

\begin{itemize}
  \item \textbf{Private Relay (الترحيل الخاص)}: "VPN" بنسخة Apple. يمر مرور Safari عبر خوادم Apple ثم يُرسل للموقع، ولا أحد (بما فيهم Apple) يستطيع رؤية من أنت وماذا زرت في آن واحد. مناسب لحماية الخصوصية على Wi-Fi عام. مسار الإعداد:\texttt{الإعدادات > اسمك > iCloud > الترحيل الخاص}.
  \item \textbf{Hide My Email (إخفاء البريد)}: عند التسجيل في المواقع/التطبيقات يولّد عنوان بريد عشوائي لمرة واحدة، تُحوّل الرسائل تلقائياً إلى بريدك الحقيقي. إذا لم تعد تريد رسائل هذا الموقع، أغلق هذا العنوان العشوائي مباشرة، ولا يُكشف بريدك الحقيقي. مسار الإعداد:\texttt{الإعدادات > اسمك > iCloud > إخفاء البريد}.
  \item \textbf{نطاق بريد إلكتروني مخصص}: يمكنك استخدام نطاقك الخاص (مثل @اسم عائلتك.com) لإرسال واستقبال بريد iCloud، مناسب للعلامة الشخصية أو إدارة بريد الأسرة.
\end{itemize}

تلميح: Private Relay غير متاح في بعض البلدان/المناطق، وقد يحتاج المستخدمون المحليون لحلول أخرى مكملة.

\end{tcbinfo}

\begin{tcbwarn}{iCloud Advanced Data Protection (ADP) — التشفير من طرف إلى طرف الحقيقي}
في نهاية 2022 أطلقت Apple\textbf{iCloud Advanced Data Protection (ADP)}ومدّت التشفير من طرف إلى طرف (E2EE) من البيانات القليلة السابقة كالصحة وكلمات المرور، إلى\textbf{معظم بيانات iCloud}— الصور والملاحظات والنسخ الاحتياطي السحابي لـ iCloud ومعلومات الجهاز وما مجموعه 23 فئة من البيانات.

\begin{itemize}
  \item \textbf{ما هو E2EE}: مفتاح التشفير موجود على جهازك فقط، حتى خوادم Apple لا تملكه. حتى لو اخترقت Apple أو طلبت إنفاذ القانون البيانات، لا تستطيع Apple تقديم بيانات يمكن فك تشفيرها.
  \item \textbf{مسار التفعيل}:\texttt{الإعدادات > اسمك > iCloud > حماية البيانات المتقدمة}، واتبع التعليمات لإعداد جهة اتصال الاستعادة/مفتاح الاستعادة (\textbf{يجب فعله}، إذا فقدت المفتاح لا يمكن استرجاع بياناتك).
  \item \textbf{تحذيرات}: (1) بعد التفعيل احرص على حفظ مفتاح الاست\_recovery بشكل آمن، وإلا لا يمكن استرجاع البيانات عند فقدان الحساب؛ (2) عدد قليل من الوظائف مثل ألبوم الصور المشترك ومكتبة صور iCloud المشتركة لا تشارك في E2EE؛ (3) Apple ID للمنطقة الصينية لا يدعم ADP بعد، ويحتاج التبديل إلى حساب منطقة أخرى.
\end{itemize}

نصيحة بكلمة واحدة: إذا كان حسابك يحتوي على بيانات خصوصية مهمة، ومنطقة الحساب تدعم، يُنصح بشدة بتفعيل ADP — هذا أقوى حماية لبياناتك السحابية يمكن أن يحصل عليها المستخدم العادي.

\end{tcbwarn}

\booksection{رابعاً: المزامنة والنسخ الاحتياطي، لا تخلط بينهما}

هاتان الكلمتان متشابهتان في الظاهر، لكن العواقب مختلفة تماماً:

\begin{itemize}
  \item \textbf{المزامنة}ثنائية الاتجاه: إذا حذفت من الهاتف، تُحذف من السحابة؛ وإذا حذفت من السحابة، يتبعها الهاتف. صور iCloud وجهات الاتصال والملاحظات كلها مزامنة. الميزة هي الاتساق بين الأجهزة المتعددة، والعيب أن "الحذف يكون نهائياً".
  \item \textbf{النسخ الاحتياطي}أحادي الاتجاه: يلتقط الهاتف "لقطة" من حالته الحالية ويخزنها في السحابة. بعد ذلك، إذا حذفت شيئاً من الهاتف، تظل النسخة في النسخة الاحتياطية السحابية. iCloud Backup والنسخ الاحتياطي المحلي على الحاسوب كلاهما نسخ احتياطي. الميزة هي أن البيانات أكثر أماناً، والعيب هو أنها تشغل مساحة ولا تتحدث في الوقت الفعلي.
\end{itemize}

\booksection{خامساً: مصفوفة حسابات مستخدم أندرويد}

الوضع في أندرويد أكثر تعقيداً قليلاً من Apple، لأن بياناتك قد تكون موزعة على عدة حسابات:

\begin{itemize}
  \item \textbf{حساب العلامة التجارية}(Huawei/Xiaomi/OPPO/vivo): يدير مساحتك السحابية، والعثور على الهاتف، ومتجر التطبيقات، والمحفظة. لكل علامة خدمتها السحابية الخاصة، عادةً بنموذج أساس مجاني + توسع مدفوع.
  \item \textbf{حساب Google}: إذا كنت تستخدم Gmail وGoogle Photos وجهات اتصال Google، فكل هذه الأشياء في حساب Google. تمنحك Google مجاناً 15 غيغابايت، أكثر كرماً من Apple، لكنها مشتركة مع Gmail وDrive.
  \item \textbf{WeChat/QQ}: مستقلة عن نظام الهاتف. سجل محادثاتك ومفضلاتك في WeChat ليست في حساب العلامة التجارية ولا في حساب Google، بل فقط على خوادم WeChat نفسها.
\end{itemize}

لذلك، عندما يقوم مستخدم أندرويد بالنسخ الاحتياطي، يجب أن يتحقق بشكل منفصل من: النسخ الاحتياطي السحابي للعلامة التجارية، والنسخ الاحتياطي من Google، والنسخ الاحتياطي لسجل محادثات WeChat — لا غنى عن أي منها.

\booksection{سادساً: المشاركة العائلية: ليستفيد أفراد العائلة أيضاً}

تتيح “المشاركة العائلية” من Apple لك ولعائلتك مشاركة مشتريات التطبيقات والموسيقى والأفلام ومساحة iCloud، بل ومشاركة الموقع. يمكن إضافة حتى 6 أشخاص. طريقة الإعداد: الإعدادات اسمك المشاركة العائلية. لدى Huawei وXiaomi ميزات مشابهة، يمكن العثور على “المشاركة العائلية” أو “السحابة العائلية” في إعدادات الحساب.

سيناريو عملي جداً: إعداد المشاركة العائلية على هاتف والديك، لتتمكن من رؤية موقعهما في أي وقت (بموافقتهما)، فإذا ضاع أحدهما أو تعذر التواصل معه، يكون لديك وسيلة لتحديد الموقع.

\begin{tcbwarn}{عواقب اختراق الحساب أخطر مما تتخيل}
إذا حصل القراصنة على كلمة مرور حسابك، يمكنه: قفل هاتفك عن بُعد وابتزازك بدفع فدية، والاطلاع على صورك وجهات اتصالك، واستخدام وسيلة الدفع الخاصة بك للشراء. الأسوأ أنه إذا كان حسابك مرتبطاً بأشياء كثيرة، فإن استعادة الحساب نفسه كابوس. لذلك، فعّل المصادقة الثنائية الآن.

\end{tcbwarn}

\begin{tcbpractice}{تطبيق عملي: أجرِ فحصاً شاملاً لحسابك}
\begin{enumerate}
  \item افتح إعدادات الهاتف، ادخل صفحة الحساب، وتأكد من تفعيل المصادقة الثنائية.
  \item تحقق من إعداد "أرقام الهاتف الموثوقة" و"مفتاح الاسترجاع"، واكتب المفتاح على ورقة.
  \item ادخل إدارة iCloud أو المساحة السحابية، لترى ما يتم مزامنته، وأوقف ما لا تحتاجه.
  \item تأكد من تفعيل "العثور على هاتفي".
  \item إذا كنت تستخدم أندرويد، تحقق من تغطية النسخ الاحتياطي السحابي للعلامة التجارية، والنسخ الاحتياطي من Google، والنسخ الاحتياطي لـ WeChat.
\end{enumerate}

\end{tcbpractice}
